\begin{resumo}
\vspace{\onelineskip}
\addcontentsline{toc}{chapter}{Resumo}
O presente trabalho discorreu sobre a importância da coleta de dados pessoais na sociedade moderna.
Descreveu-se por quais métodos ela ocorre, em que circunstâncias, e, principalmente, como as quantidades massivas de dados coletados fundamentam o enorme poder político e econômico das assim chamadas Big Techs.
Analisou-se os algoritmos de aprendizado de máquina, que processam a maioria desses dados, quanto ao seu funcionamento, suas limitações, sua transparência, sua necessidade de tantos dados e sua capacidade de dar significado a observações aparentemente banais.
Caracterizou-se o Vale do Silício como uma potência política e ideológica que, fundamentada na coleta de dados, pode utilizar-se do neoliberalismo dos anos 80 e 90 para formar gigantescos monopólios pelo meio digital.
Por fim, apresentou-se um modo de produção alternativo: o \textit{software} livre, que, ao permitir total publicidade dos seus mecanismos de funcionamento, pode ser uma solução para os problemas apresentados, possibilitando uma máxima transparência na coleta e processamento de dados.

\textbf{Palavras-chave:} dados pessoais; privacidade; transparência; vigilância; Vale do Silício.
\end{resumo}

\renewcommand{\resumoname}{Abstract}
\begin{resumo}
\begin{otherlanguage}{english}
\vspace{\onelineskip}
\addcontentsline{toc}{chapter}{Abstract}
In this monograph, we analyse the role of personal data collection in modern society, by describing the technical methods through which it works, and the political and economic circumstances that surround it.
We show how data collection enabled the rise of major monopolist organizations with gigantic political and economic power: the so-called Big Techs.
By understanding the inner workings of the algorithms that process these data through the lenses of process transparency, we describe their main constitutional characteristics, limitations, why they require such large amounts of data, and their capacities of signifying seemingly unimportant observations.
We characterize the Silicon Valley as a political and ideological force that, based on the mechanisms of personal data collection, could profit from late twentieth century neoliberalism and form monopolies around the digital space.
At last, we present free software as an alternative system of software production that can enable full transparency and, by such, be a potential solution for the problems regarding personal data collection.

\textbf{Keywords:} personal data; privacy; transparency; surveillance; Silicon Valley.
\end{otherlanguage}
\end{resumo}
